% Pattern: head/cite. A cite by number after the number or the name of a
% head, as "Lemma 1[6,9]", is the source, and so is a cite after a named
% conjecture, as "Vizing's conjecture \cite{key}.". The words after the
% first word of a named theorem can be in lower case. Seen in the gold pilot
% (set X).
% want head-1 statement/lemma 1
% want head-2 statement/theorem -
% want head-3 conjecture -
% want head-4 statement/theorem -
\documentclass{article}
\begin{document}
\section{One}
We recall the results that we use.

\noindent{\bf Lemma 1[6,9]} {\it Every compact space is bounded.}

{\bf Theorem (Heine's Theorem) [9]}. Every closed bounded set is compact.

\noindent\textbf{Bounded conjecture \cite{key}.} Every space is bounded.

\textbf{Analytic continuous fundamental theorem.} Every finite set is closed.
\end{document}
